\problemname{The Five Card Game}
Alf and Beata were two youngsters who lived a long, long time ago, in the days before one could spend one's afternoons on programming competitions.
Their lives were thus much more boring than the lives of today's youth.
How can one survive without computers, you might ask.
The answer is very simple: you play cards!

Our two youngsters loved playing cards, and often had large piles of card decks that they received as Christmas gifts every year.
To avoid filling their rooms with card decks, Beata challenged her friend to a game every evening --- five card.

Five card is played by two players (in our case, Alf and Beata), with a standard deck of $52$ cards where each player has been dealt five cards.
The players will now play five so-called tricks.
A trick works as follows: a player $A$ chooses one of their cards and plays it on the table.
The opponent then chooses one of their cards (which must be of the same suit as the opponent's, if they have one) and plays it on the table.
If the opponent did not have a card of the same suit, player $A$ won the trick, otherwise the player whose card had the highest value won.
In the next trick, the player who won the previous trick starts by playing a card.
They take turns in this way until both players are out of cards, i.e.\ after $5$ tricks.

Alf starts by playing a card in the first trick.
Given the order in which the two players played their cards, can you determine who won the last trick, and whether either (or both) of the players cheated?

\section*{Input}
The input begins with five lines containing Alf's cards, in the order the cards were played.
Then follow five lines with Beata's cards, in the order the cards were played. All cards are distinct.

A card has the form \texttt{suit value}, where the suit is one of \texttt{R, K, H, S} (for diamonds, clubs, hearts, spades), and the value is an integer between 2 and 14.

\section*{Output}
Start by outputting the player who won the last trick: \texttt{A} if Alf won, and \texttt{B} if Beata won.

If any player cheated, output on the next line the players who cheated.
If Alf cheated, output \texttt{A fuskar},
if Beata cheated, output \texttt{B fuskar},
and if both cheated, output \texttt{A och B fuskar}.

\section*{Scoring}
Your solution will be tested on a set of test groups. To earn points for a group, you must pass all test cases in that group.

\noindent
\begin{tabular}{| l | l | p{12cm} |}
  \hline
  \textbf{Group} & \textbf{Points} & \textbf{Constraints} \\ \hline
  $1$    & $30$      & All cards played have the same suit \\ \hline
  $2$    & $30$      & No player cheated \\ \hline
  $3$    & $40$      & No additional constraints.  \\ \hline
\end{tabular}
